\chapter{Исследовательская часть}

\section{Технические характеристики}

Технические характеристики устройства,на котором проводились замеры:

\begin{itemize}
	\item AMD Ryzen 7 5800U (16) @ 4.51 ГГц;
	\item Оперативная память: 16 Гб;
	\item Ядро ОС Linux 6.10.10-arch1-1.
\end{itemize}

При проведении замеров ноутбук был включен в сеть и были запущены только системные приложения и реализованные программы.

\section{Временные характеристики}
В таблице \ref{tbl:time} приведены временные характеристики, полученные в результате замеров времени алгоритмов. Для каждой длины замер проводился 100 раз, при этом для каждой итерации создавались 2 случайные строки заданной длины и использовались как аргументы к алгоритмам. Итоговое время взято как среднее арифметическое всех полученных замеров.

\begin{table}[H]
	\small
	\begin{center}
		\begin{threeparttable}
			\caption{Временные характеристики}
			\begin{tabular}{|c|r|r|}
				\hline
				Длина строки, символы & Рекурсивный алгоритм, нс & Алгоритм с кешем, нс \\
				\hline
				2 & 2 915.99  & 2 227.65  \\
				\hline
				3 & 11 837.55  & 3 451.7  \\
				\hline
				4 & 55 621.85  & 4 863.58  \\
				\hline
				5 & 257 908.21  & 6 819.9  \\
				\hline
				6 & 1 375 881.37  & 9 031.91  \\
				\hline
				7 & 6 948 788.65  & 11 651.07  \\
				\hline
				8 & 35 773 610.35  & 14 515.6  \\
				\hline
			\end{tabular}
                \label{tbl:time}
		\end{threeparttable}
	\end{center}
\end{table}

Полученные замеры также можно увидеть на графике \ref{fig:time}:

\begin{figure}[H]
	\centering
	\includegraphics[width=1\textwidth]{time}
	\caption{График зависимости времени работы алгоритмов от длины строк}
	\label{fig:time}
\end{figure}

Как видно из таблицы \ref{tbl:time} и графика \ref{fig:time}, рекурсивный алгоритм существенно уступает итерационному по временным характеристикам.

\pagebreak

\section{Емкостные характеристики}

В таблице \ref{tbl:mem} приведены временные характеристики, полученные в результате замеров памяти алгоритмов с помощью модуля tracemalloc. Так как для рекурсивной реализации размер выделяемой памяти существенно зависит от глубины рекурсии, то есть от входных строк, то для каждой длины строки для рекурсивного алгоритма алгоритм был проведен 8 раз на случайных строках. Для итерационного алгоритма выделяемая память не зависит от входных строк(при одинаковой длине), так как размер матрицы кеша зависит только от размеров строк. За характеристику памяти взята пиковая память, то есть максимальный объем памяти выделенный функции во время её выполнения.

\begin{table}[H]
	\small
	\begin{center}
		\begin{threeparttable}
			\caption{Емкостные характеристики}
			\begin{tabular}{|c|r|r|}
				\hline
				Длина строки, символы & Рекурсивный алгоритм, байты & Алгоритм с кешем, байты \\
				\hline
				2 & 48  & 232  \\
				\hline
				3 & 247  & 288  \\
				\hline
				4 & 288  & 392  \\
				\hline
				5 & 312  & 480  \\
				\hline
				6 & 404  & 584  \\
				\hline
				7 & 498  & 704  \\
				\hline
				8 & 674  & 904  \\
				\hline
			\end{tabular}
                \label{tbl:mem}
		\end{threeparttable}
	\end{center}
\end{table}

Также зависимость из таблицы отражена на графике \ref{fig:size1}.

\begin{figure}[H]
	\centering
	\includegraphics[width=0.9\textwidth]{size1}
	\caption{График зависимости требуемой памяти для работы алгоритмов от длины входных строк}
	\label{fig:size1}
\end{figure}

Как видно из графика \ref{fig:size1} и таблицы \ref{tbl:mem}, при небольших размерах строки рекурсивная реализация выигрывает итерационную, однако начиная с длин строк 4, вариант с кешем становится эффективнее, так как реализация рекурсии требует больший объем памяти под входные строки и вызовы.

\pagebreak

\section{Вывод}

В данном разделе было проведено исследование временных и ресурсных характеристик алгоритмов поиска расстояния Левенштейна: рекурсивного и итерационного с кешированием.

Из исследования можно сделать вывод, что по времени выполнения рекурсивный алгоритм значительно уступает итерационному. Однако, по использованию ресурсов, при длине входных строк $\leq 3$ рекурсивный подход оказывается более эффективным. С увеличением длины строк до $\geq 4$ итерационный метод становится предпочтительнее, поскольку рекурсивный алгоритм вызывает множество функций, что приводит к увеличению вычислительных затрат и нагрузки на стек вызовов.

\clearpage
